\textbf{Phần 1. Trắc nghiệm lựa chọn}
\setcounterpageref{ex}{0}

\begin{ex} %[1D5V18]
Tính giá trị của biểu thức: $3 + 7 - 4$.
\choice
{$5$}
{\True $6$}
{$7$}
{$4$}
\loigiai{
Thực hiện từ trái sang phải: $3 + 7 = 10$; sau đó $10 - 4 = 6$.
\textbf{Đáp số:} $6$.\par
\textbf{Đáp án của câu hỏi là B}
}
\end{ex}

\begin{ex} %[1D5V18]
Tính giá trị của biểu thức: $10 - 5 + 3$.
\choice
{\True $8$}
{$7$}
{$9$}
{$6$}
\loigiai{
Thực hiện từ trái sang phải: $10 - 5 = 5$; sau đó $5 + 3 = 8$.
\textbf{Đáp số:} $8$.\par
\textbf{Đáp án của câu hỏi là A}
}
\end{ex}

\begin{ex} %[1D5V18]
Tính giá trị của biểu thức: $9 - 2 - 4 + 5$.
\choice
{$7$}
{$9$}
{\True $8$}
{$6$}
\loigiai{
Thực hiện lần lượt: $9 - 2 = 7$; $7 - 4 = 3$; $3 + 5 = 8$.
\textbf{Đáp số:} $8$.\par
\textbf{Đáp án của câu hỏi là C}
}
\end{ex}

\begin{ex} %[1D5V18]
Điền số thích hợp vào ô trống: $4 + \dots = 10 - 1$.
\choice
{$4$}
{$6$}
{$3$}
{\True $5$}
\loigiai{
Vế phải: $10 - 1 = 9$.\nTa có: $4 + \dots = 9 \Rightarrow$ số cần điền là: $9 - 4 = 5$.
\textbf{Đáp số:} $5$.\par
\textbf{Đáp án của câu hỏi là D}
}
\end{ex}

\begin{ex} %[1D5V18]
Điền số thích hợp vào ô trống: $8 - \dots = 2 + 3$.
\choice
{$2$}
{\True $3$}
{$4$}
{$1$}
\loigiai{
Vế phải: $2 + 3 = 5$.\nTa có: $8 - \dots = 5 \Rightarrow$ số cần điền là: $8 - 5 = 3$.
\textbf{Đáp số:} $3$.\par
\textbf{Đáp án của câu hỏi là B}
}
\end{ex}

\begin{ex} %[1D5V18]
Tìm số tự nhiên $x$ thỏa mãn: $x + 3 = 5 + 2$.
\choice
{\True $4$}
{$3$}
{$5$}
{$2$}
\loigiai{
Vế phải: $5 + 2 = 7$.\nTa có: $x + 3 = 7 \Rightarrow x = 7 - 3 = 4$.
\textbf{Đáp số:} $4$.\par
\textbf{Đáp án của câu hỏi là A}
}
\end{ex}

\begin{ex} %[1D5V18]
Tìm số tự nhiên $y$ thỏa mãn: $10 - y = 4 + 3$.
\choice
{$2$}
{$4$}
{\True $3$}
{$1$}
\loigiai{
Vế phải: $4 + 3 = 7$.\nTa có: $10 - y = 7 \Rightarrow y = 10 - 7 = 3$.
\textbf{Đáp số:} $3$.\par
\textbf{Đáp án của câu hỏi là C}
}
\end{ex}

\begin{ex} %[1D5V18]
Lan có 4 cái kẹo, mẹ cho Lan thêm 5 cái kẹo nữa. Sau đó Lan cho bạn 3 cái kẹo. Hỏi Lan còn lại bao nhiêu cái kẹo?
\choice
{$5$}
{$7$}
{$4$}
{\True $6$}
\loigiai{
Số kẹo Lan có sau khi mẹ cho là: $4 + 5 = 9$ (cái kẹo).\nSố kẹo Lan còn lại là: $9 - 3 = 6$ (cái kẹo).
\textbf{Đáp số:} $6$.\par
\textbf{Đáp án của câu hỏi là D}
}
\end{ex}

\begin{ex} %[1D5V18]
Trên đĩa có 10 quả dâu tây. Nam ăn 4 quả, Mai ăn 3 quả. Hỏi trên đĩa còn lại bao nhiêu quả dâu tây?
\choice
{$2$}
{\True $3$}
{$4$}
{$1$}
\loigiai{
Tổng số dâu tây hai bạn đã ăn là: $4 + 3 = 7$ (quả).\nSố dâu tây còn lại trên đĩa là: $10 - 7 = 3$ (quả).
\textbf{Đáp số:} $3$.\par
\textbf{Đáp án của câu hỏi là B}
}
\end{ex}

\begin{ex} %[1D5V18]
Bác nông dân nuôi 5 con bò sữa và 4 con bò thịt. Bác bán đi 2 con bò thịt. Hỏi đàn bò nhà bác còn lại bao nhiêu con cả hai loại?
\choice
{\True $7$}
{$6$}
{$8$}
{$5$}
\loigiai{
Số bò thịt còn lại là: $4 - 2 = 2$ (con).\nTổng số bò còn lại là: $5 + 2 = 7$ (con bò).
\textbf{Đáp số:} $7$.\par
\textbf{Đáp án của câu hỏi là A}
}
\end{ex}

\begin{ex} %[1D5V18]
Tìm một số biết rằng khi lấy số đó cộng với 3 rồi trừ đi 1 thì được số lớn nhất có một chữ số.
\choice
{$6$}
{$8$}
{\True $7$}
{$5$}
\loigiai{
Số lớn nhất có một chữ số là 9.\nTa có: $x + 3 - 1 = 9 \Rightarrow x + 2 = 9 \Rightarrow x = 9 - 2 = 7$.
\textbf{Đáp số:} $7$.\par
\textbf{Đáp án của câu hỏi là C}
}
\end{ex}

\begin{ex} %[1D5V18]
Tìm một số biết rằng lấy 10 trừ đi số đó rồi cộng thêm 4 thì được 8.
\choice
{$5$}
{$7$}
{$4$}
{\True $6$}
\loigiai{
Ta có: $10 - x = 8 - 4 = 4 \Rightarrow x = 10 - 4 = 6$.
\textbf{Đáp số:} $6$.\par
\textbf{Đáp án của câu hỏi là D}
}
\end{ex}

\textbf{Phần 2. Trắc nghiệm đúng sai}
\setcounterpageref{ex}{0}

\begin{ex} %[1D5V18]
Xét tính đúng hoặc sai của các khẳng định toán học sau đây:
\choiceTF
{Hùng có 8 viên bi, Dũng có ít hơn Hùng 3 viên bi. cả hai bạn có $12$ viên bi.}
{\True Có hai thùng đựng nước: thùng thứ nhất đựng 6 lít, thùng thứ hai đựng 4 lít. Người ta dùng hết 5 lít nước từ hai thùng. cả hai thùng còn lại $5$ lít nước.}
{\True Số thích hợp cần điền vào $2 + 2 + 2 + \dots = 10$ là $4$.}
{Số thích hợp cần điền vào $10 - 3 - 3 - \dots = 1$ là $2$.}
\loigiai{
\begin{itemize}
    \item Ý a) Sai: Kết quả đúng là $13$, không phải $12$. Số bi của Dũng là: $8 - 3 = 5$ (viên bi).\nCả hai bạn có tất cả: $8 + 5 = 13$ (viên bi).
\textbf{Đáp số:} $13$.
    \item Ý b) Đúng: Tổng số nước ban đầu ở cả hai thùng là: $6 + 4 = 10$ (lít).\nSố lít nước còn lại là: $10 - 5 = 5$ (lít).
\textbf{Đáp số:} $5$.
    \item Ý c) Đúng: Ta có: $2 + 2 + 2 = 6$.\nSố cần điền là: $10 - 6 = 4$.
\textbf{Đáp số:} $4$.
    \item Ý d) Sai: Kết quả đúng là $3$, không phải $2$. Ta có: $10 - 3 - 3 = 4$.\nSố cần điền là: $4 - 1 = 3$.
\textbf{Đáp số:} $3$.
\end{itemize}
\begin{itemize}
    \item a) \textbf{Sai}.
    \item b) \textbf{Đúng}.
    \item c) \textbf{Đúng}.
    \item d) \textbf{Sai}.
\end{itemize}
}
\end{ex}

\begin{ex} %[1D5V18]
Xét tính đúng hoặc sai của các bài toán và phép tính sau đây:
\choiceTF
{\True Kết quả của bài toán «Tìm hai số có tổng bằng 10 và hiệu bằng 2. Số bé hơn là số mấy» là $4$.}
{Kết quả của bài toán «Tìm hai số có tổng bằng 9 và hiệu bằng 5. Số lớn hơn là số mấy» là $6$.}
{\True Một đội đồng diễn gồm 10 bạn xếp thành hai hàng. Hàng trước có 6 bạn. hàng sau có $4$ bạn.}
{Số thích hợp cần điền vào $2 + 2 + 2 + \dots = 10$ là $5$.}
\loigiai{
\begin{itemize}
    \item Ý a) Đúng: Hai số có tổng bằng 10 và hiệu bằng 2 là 6 và 4 ($6 + 4 = 10$ và $6 - 4 = 2$).\nSố bé hơn là số $4$.
\textbf{Đáp số:} $4$.
    \item Ý b) Sai: Kết quả đúng là $7$, không phải $6$. Hai số có tổng bằng 9 và hiệu bằng 5 là 7 và 2 ($7 + 2 = 9$ và $7 - 2 = 5$).\nSố lớn hơn là số $7$.
\textbf{Đáp số:} $7$.
    \item Ý c) Đúng: Số bạn ở hàng sau là: $10 - 6 = 4$ (bạn).
\textbf{Đáp số:} $4$.
    \item Ý d) Sai: Kết quả đúng là $4$, không phải $5$. Ta có: $2 + 2 + 2 = 6$.\nSố cần điền là: $10 - 6 = 4$.
\textbf{Đáp số:} $4$.
\end{itemize}
\begin{itemize}
    \item a) \textbf{Đúng}.
    \item b) \textbf{Sai}.
    \item c) \textbf{Đúng}.
    \item d) \textbf{Sai}.
\end{itemize}
}
\end{ex}

\textbf{Phần 3. Trả lời ngắn}
\setcounterpageref{ex}{0}

\begin{ex} %[1D5V18]
Hàng trước nhiều hơn hàng sau bao nhiêu bạn học sinh?
\shortans{2}
\loigiai{
\textbf{Đáp số:} $2$\par
Số bạn hàng trước nhiều hơn hàng sau là: $6 - 4 = 2$ (bạn).
\textbf{Đáp số:} $2$.
}
\end{ex}

\begin{ex} %[1D5V18]
Tìm số tự nhiên $m$ sao cho: $m + m + 1 = 9$. Giá trị của $m$ là bao nhiêu?
\shortans{4}
\loigiai{
\textbf{Đáp số:} $4$\par
Ta có: $m + m = 9 - 1 = 8 \Rightarrow m = 4$ (vì $4 + 4 = 8$).
\textbf{Đáp số:} $4$.
}
\end{ex}

\begin{ex} %[1D5V18]
Tìm số tự nhiên $n$ sao cho: $10 - n - n = 4$. Giá trị của $n$ là bao nhiêu?
\shortans{3}
\loigiai{
\textbf{Đáp số:} $3$\par
Ta có: $n + n = 10 - 4 = 6 \Rightarrow n = 3$ (vì $3 + 3 = 6$).
\textbf{Đáp số:} $3$.
}
\end{ex}

\begin{ex} %[1D5V18]
Một hàng rào có 4 cây cọc. Mỗi khoảng cách giữa hai cây cọc dài 2 mét. Hỏi hàng rào đó dài tất cả bao nhiêu mét?
\shortans{6}
\loigiai{
\textbf{Đáp số:} $6$\par
4 cây cọc có 3 khoảng cách.\nChiều dài hàng rào là: $2 + 2 + 2 = 6$ (mét).
\textbf{Đáp số:} $6$.
}
\end{ex}

\begin{ex} %[1D5V18]
Một cửa hàng buổi sáng bán được 5 chiếc quạt, buổi chiều bán được 4 chiếc quạt. Người ta nhập về thêm 3 chiếc quạt mới. Hỏi số quạt bán được trong cả ngày là bao nhiêu chiếc?
\shortans{9}
\loigiai{
\textbf{Đáp số:} $9$\par
Số quạt bán được trong cả ngày là: $5 + 4 = 9$ (chiếc quạt).
\textbf{Đáp số:} $9$.
}
\end{ex}

\begin{ex} %[1D5V18]
Tính giá trị của biểu thức: $10 - 1 - 2 - 3 - 4$.
\shortans{0}
\loigiai{
\textbf{Đáp số:} $0$\par
Thực hiện trừ lần lượt: $10 - 1 = 9$; $9 - 2 = 7$; $7 - 3 = 4$; $4 - 4 = 0$.
\textbf{Đáp số:} $0$.
}
\end{ex}
